\documentclass[11pt,letterpaper]{article}
\usepackage[technical-reference]{cornell-notes}
\usepackage{needspace}

\title{Clause 23: Strings Library - Cornell Notes}
\author{}
\date{}
\setDocTitle{Clause 23: Strings Library}
\setDocSubtitle{C++ 2024 Cornell Notes}
\setDocOwner{C++ 2024 Cornell Notes}
\setDocAuthor{}
\setDocDate{}
\setCornellCollection{C++ 2024 Cornell Notes}
\setCornellUnitType{Clause}
\setCornellUnitNumber{23}
\setCornellUnitTitle{Strings Library}
\setCornellCompanionLabel{C++ Technical Reference}
\hypersetup{pdftitle={Clause 23: Strings Library - Cornell Notes},pdfsubject={C++ technical study notes},pdfauthor={}}

\begin{document}
\maketitle
\begin{CornellReferenceOrientationBox}
\textbf{Purpose.} Defines character traits, owning strings, string views, null-terminated sequence utilities, numeric conversions, and related hashing and I/O. It distinguishes ownership, bounds, encoding units, and invalidation behavior.
\par\textbf{Role.} The primary character-sequence abstraction layer.
\par\textbf{Principal dependencies.} Uses the library-wide rules of Clause 16 together with the applicable core-language concepts and supporting library clauses.
\par\textbf{Primary audience.} programmers and text-processing library authors.
\par\textbf{Major subject groups.} General; Character traits; String view classes; String classes; Null-terminated sequence utilities.
\end{CornellReferenceOrientationBox}
\section{Learning Objectives}
\begin{itemize}[label=\textcolor{CornellReferenceTeal}{$\blacktriangleright$}]
\item Define the governing ideas of General and apply its stated contract at a call site.
\item Identify the governing ideas of Character traits and apply its stated contract at a call site.
\item Distinguish the governing ideas of String view classes and apply its stated contract at a call site.
\item Explain the governing ideas of String classes and apply its stated contract at a call site.
\item Trace the governing ideas of Null-terminated sequence utilities and apply its stated contract at a call site.
\item Apply the governing ideas of General and apply its stated contract at a call site.
\item Analyze the governing ideas of Character traits requirements and apply its stated contract at a call site.
\item Evaluate the governing ideas of Traits typedefs and apply its stated contract at a call site.
\item Diagnose the governing ideas of char\_\allowbreak{}traits specializations and apply its stated contract at a call site.
\item Compare the governing ideas of General and apply its stated contract at a call site.
\item Verify the governing ideas of Header <string\_\allowbreak{}view> synopsis and apply its stated contract at a call site.
\item Relate the governing ideas of Class template basic\_\allowbreak{}string\_\allowbreak{}view and apply its stated contract at a call site.
\item Classify the governing ideas of Non-member comparison functions and apply its stated contract at a call site.
\item Justify the governing ideas of Inserters and extractors and apply its stated contract at a call site.
\item Audit the governing ideas of Hash support and apply its stated contract at a call site.
\item Summarize the governing ideas of Suffix for basic\_\allowbreak{}string\_\allowbreak{}view literals and apply its stated contract at a call site.
\end{itemize}
\section{Normative Structure Map}
\small\begin{longtable}{@{}>{\RaggedRight\arraybackslash}p{0.13\textwidth}>{\RaggedRight\arraybackslash}p{0.27\textwidth}>{\RaggedRight\arraybackslash}p{0.19\textwidth}>{\RaggedRight\arraybackslash}p{0.31\textwidth}@{}}
\toprule\textbf{Subclause} & \textbf{Subject} & \textbf{Rule category} & \textbf{Key dependencies / entities}\\\midrule\endfirsthead
\toprule\textbf{Subclause} & \textbf{Subject} & \textbf{Rule category} & \textbf{Key dependencies / entities}\\\midrule\endhead
\bottomrule\endfoot
23.1 & General & library semantics & Uses the library-wide rules of Clause 16 together with the applicable core-language concepts and supporting library clauses.\\
23.2 & Character traits & library semantics & General, Character traits requirements, Traits typedefs, char\_\allowbreak{}traits specializations\\
23.3 & String view classes & library semantics & General, Header <string\_\allowbreak{}view> synopsis, Class template basic\_\allowbreak{}string\_\allowbreak{}view, Non-member comparison functions, and related subclauses\\
23.4 & String classes & library semantics & General, Header <string> synopsis, Class template basic\_\allowbreak{}string, Non-member functions, and related subclauses\\
23.5 & Null-terminated sequence utilities & library semantics & Header <cctype> synopsis, Header <cwctype> synopsis, Header <cstring> synopsis, Header <cwchar> synopsis, and related subclauses\\
\end{longtable}\normalsize
\begin{CornellReferenceNormativeBox}A heading maps the clause; it does not replace the detailed conditions. For each operation, read requirements, effects, result, exceptions, complexity, remarks, and notes with their stated normative force.\end{CornellReferenceNormativeBox}
\subsection*{Rule Classification Guide}
\small\begin{longtable}{@{}>{\RaggedRight\arraybackslash}p{0.25\textwidth}>{\RaggedRight\arraybackslash}p{0.67\textwidth}@{}}\toprule\textbf{Label or category} & \textbf{How to read it}\\\midrule\endfirsthead\toprule\textbf{Label or category} & \textbf{How to read it (continued)}\\\midrule\endhead\bottomrule\endfoot
Constraints & Conditions controlling whether a declaration, overload, or specialization participates in the relevant context.\\
Mandates & Requirements checked for the described declaration or specialization; failure has the stated well-formedness consequence.\\
Preconditions & Obligations the caller establishes before evaluation; violating one forfeits the operation's contract.\\
Effects / Returns / Postconditions & Separate the state change, produced result, and state guaranteed after successful completion.\\
Throws / Error conditions & State how failure is reported and which exceptional paths are permitted or required.\\
Complexity & Bounds or exact counts for the operations named by the specification.\\
Remarks / Recommended practice & Remarks refine applicability or participation; recommended practice guides implementations without silently becoming a program requirement.\\
Exposition-only & A device for describing semantics, not a public name on which a program can depend.\\
\end{longtable}\normalsize
\section{Cornell Cue-and-Notes}
\Needspace{8\baselineskip}
\subsection*{23.1 General}
\begin{longtable}{@{}>{\RaggedRight\arraybackslash\columncolor{CornellReferenceCueGray}}p{0.28\textwidth}>{\RaggedRight\arraybackslash}p{0.64\textwidth}@{}}
\rowcolor{CornellReferenceNavy}\color{white}\textbf{Cue / recall prompt} & \color{white}\textbf{Notes}\\\endfirsthead
\rowcolor{CornellReferenceNavy}\color{white}\textbf{Cue / recall prompt} & \color{white}\textbf{Notes (continued)}\\\endhead
\arrayrulecolor{CornellReferenceLineGray}
\textbf{What contract governs General?}\\[-1mm]{\scriptsize\CornellClauseRef{23.1}} & Frames the terminology, applicability, and relationships needed to interpret General; detailed rules remain in the following subclauses.\\\hline
\end{longtable}
\begin{CornellReferenceSummaryBox}General supplies the library semantics portion of this clause. The key review move is to connect its local wording to the clause-wide model, then classify each condition by normative force before relying on it.\end{CornellReferenceSummaryBox}
\Needspace{8\baselineskip}
\subsection*{23.2 Character traits}
\begin{longtable}{@{}>{\RaggedRight\arraybackslash\columncolor{CornellReferenceCueGray}}p{0.28\textwidth}>{\RaggedRight\arraybackslash}p{0.64\textwidth}@{}}
\rowcolor{CornellReferenceNavy}\color{white}\textbf{Cue / recall prompt} & \color{white}\textbf{Notes}\\\endfirsthead
\rowcolor{CornellReferenceNavy}\color{white}\textbf{Cue / recall prompt} & \color{white}\textbf{Notes (continued)}\\\endhead
\arrayrulecolor{CornellReferenceLineGray}
\textbf{What contract governs General?}\\[-1mm]{\scriptsize\CornellClauseRef{23.2.1}} & Frames the terminology, applicability, and relationships needed to interpret General; detailed rules remain in the following subclauses.\\\hline
\textbf{Which obligations govern Character traits requirements?}\\[-1mm]{\scriptsize\CornellClauseRef{23.2.2}} & States the admissibility and semantic obligations for Character traits requirements. Separate compile-time participation rules from caller preconditions and runtime effects.\\\hline
\textbf{What contract governs Traits typedefs?}\\[-1mm]{\scriptsize\CornellClauseRef{23.2.3}} & Specifies the facility family Traits typedefs. Read its declarations together with mandates, preconditions, effects, results, exceptions, complexity, and lifetime rules.\\\hline
\textbf{What contract governs char\_\allowbreak{}traits specializations?}\\[-1mm]{\scriptsize\CornellClauseRef{23.2.4}} & Specifies the facility family char\_\allowbreak{}traits specializations. Read its declarations together with mandates, preconditions, effects, results, exceptions, complexity, and lifetime rules.\\\hline
\end{longtable}
\begin{CornellReferenceSummaryBox}Character traits supplies the library semantics portion of this clause. The key review move is to connect its local wording to the clause-wide model, then classify each condition by normative force before relying on it.\end{CornellReferenceSummaryBox}
\Needspace{5\baselineskip}\paragraph{Detailed coverage ledger.}
\begin{description}[style=nextline,leftmargin=2.8cm,font=\normalfont\bfseries\color{CornellReferenceTeal}]
\item[23.2.4.1] \textbf{General.} Frames the terminology, applicability, and relationships needed to interpret General; detailed rules remain in the following subclauses.
\item[23.2.4.2] \textbf{struct char\_\allowbreak{}traits<char>.} Specifies the facility family struct char\_\allowbreak{}traits<char>. Read its declarations together with mandates, preconditions, effects, results, exceptions, complexity, and lifetime rules.
\item[23.2.4.3] \textbf{struct char\_\allowbreak{}traits<char8\_\allowbreak{}t>.} Specifies the facility family struct char\_\allowbreak{}traits<char8\_\allowbreak{}t>. Read its declarations together with mandates, preconditions, effects, results, exceptions, complexity, and lifetime rules.
\item[23.2.4.4] \textbf{struct char\_\allowbreak{}traits<char16\_\allowbreak{}t>.} Specifies the facility family struct char\_\allowbreak{}traits<char16\_\allowbreak{}t>. Read its declarations together with mandates, preconditions, effects, results, exceptions, complexity, and lifetime rules.
\item[23.2.4.5] \textbf{struct char\_\allowbreak{}traits<char32\_\allowbreak{}t>.} Specifies the facility family struct char\_\allowbreak{}traits<char32\_\allowbreak{}t>. Read its declarations together with mandates, preconditions, effects, results, exceptions, complexity, and lifetime rules.
\item[23.2.4.6] \textbf{struct char\_\allowbreak{}traits<wchar\_\allowbreak{}t>.} Specifies the facility family struct char\_\allowbreak{}traits<wchar\_\allowbreak{}t>. Read its declarations together with mandates, preconditions, effects, results, exceptions, complexity, and lifetime rules.
\end{description}
\Needspace{8\baselineskip}
\subsection*{23.3 String view classes}
\begin{longtable}{@{}>{\RaggedRight\arraybackslash\columncolor{CornellReferenceCueGray}}p{0.28\textwidth}>{\RaggedRight\arraybackslash}p{0.64\textwidth}@{}}
\rowcolor{CornellReferenceNavy}\color{white}\textbf{Cue / recall prompt} & \color{white}\textbf{Notes}\\\endfirsthead
\rowcolor{CornellReferenceNavy}\color{white}\textbf{Cue / recall prompt} & \color{white}\textbf{Notes (continued)}\\\endhead
\arrayrulecolor{CornellReferenceLineGray}
\textbf{What contract governs General?}\\[-1mm]{\scriptsize\CornellClauseRef{23.3.1}} & Frames the terminology, applicability, and relationships needed to interpret General; detailed rules remain in the following subclauses.\\\hline
\textbf{What contract governs Header <string\_\allowbreak{}view> synopsis?}\\[-1mm]{\scriptsize\CornellClauseRef{23.3.2}} & Catalogs the declarations made available by Header <string\_\allowbreak{}view> synopsis. The synopsis is an interface map; the detailed specifications supply constraints and behavior.\\\hline
\textbf{What contract governs Class template basic\_\allowbreak{}string\_\allowbreak{}view?}\\[-1mm]{\scriptsize\CornellClauseRef{23.3.3}} & Explains the traversal or view contract for Class template basic\_\allowbreak{}string\_\allowbreak{}view; prove domain validity, lifetime, sentinel compatibility, and any borrowing or invalidation property.\\\hline
\textbf{What contract governs Non-member comparison functions?}\\[-1mm]{\scriptsize\CornellClauseRef{23.3.4}} & Specifies the facility family Non-member comparison functions. Read its declarations together with mandates, preconditions, effects, results, exceptions, complexity, and lifetime rules.\\\hline
\textbf{What contract governs Inserters and extractors?}\\[-1mm]{\scriptsize\CornellClauseRef{23.3.5}} & Specifies the facility family Inserters and extractors. Read its declarations together with mandates, preconditions, effects, results, exceptions, complexity, and lifetime rules.\\\hline
\textbf{What contract governs Hash support?}\\[-1mm]{\scriptsize\CornellClauseRef{23.3.6}} & Specifies the facility family Hash support. Read its declarations together with mandates, preconditions, effects, results, exceptions, complexity, and lifetime rules.\\\hline
\textbf{What contract governs Suffix for basic\_\allowbreak{}string\_\allowbreak{}view literals?}\\[-1mm]{\scriptsize\CornellClauseRef{23.3.7}} & Explains the traversal or view contract for Suffix for basic\_\allowbreak{}string\_\allowbreak{}view literals; prove domain validity, lifetime, sentinel compatibility, and any borrowing or invalidation property.\\\hline
\end{longtable}
\begin{CornellReferenceSummaryBox}String view classes supplies the library semantics portion of this clause. The key review move is to connect its local wording to the clause-wide model, then classify each condition by normative force before relying on it.\end{CornellReferenceSummaryBox}
\Needspace{5\baselineskip}\paragraph{Detailed coverage ledger.}
\begin{description}[style=nextline,leftmargin=2.8cm,font=\normalfont\bfseries\color{CornellReferenceTeal}]
\item[23.3.3.1] \textbf{General.} Frames the terminology, applicability, and relationships needed to interpret General; detailed rules remain in the following subclauses.
\item[23.3.3.2] \textbf{Construction and assignment.} Governs state transfer or exchange for Construction and assignment; check self-operation, validity after movement, exception behavior, and invalidation.
\item[23.3.3.3] \textbf{Deduction guides.} Specifies the facility family Deduction guides. Read its declarations together with mandates, preconditions, effects, results, exceptions, complexity, and lifetime rules.
\item[23.3.3.4] \textbf{Iterator support.} Explains the traversal or view contract for Iterator support; prove domain validity, lifetime, sentinel compatibility, and any borrowing or invalidation property.
\item[23.3.3.5] \textbf{Capacity.} Specifies the facility family Capacity. Read its declarations together with mandates, preconditions, effects, results, exceptions, complexity, and lifetime rules.
\item[23.3.3.6] \textbf{Element access.} Specifies the facility family Element access. Read its declarations together with mandates, preconditions, effects, results, exceptions, complexity, and lifetime rules.
\item[23.3.3.7] \textbf{Modifiers.} Specifies the facility family Modifiers. Read its declarations together with mandates, preconditions, effects, results, exceptions, complexity, and lifetime rules.
\item[23.3.3.8] \textbf{String operations.} Defines the observable result of String operations together with applicable iterator-domain, ordering, complexity, invalidation, and exception conditions.
\item[23.3.3.9] \textbf{Searching.} Specifies the facility family Searching. Read its declarations together with mandates, preconditions, effects, results, exceptions, complexity, and lifetime rules.
\end{description}
\Needspace{8\baselineskip}
\subsection*{23.4 String classes}
\begin{longtable}{@{}>{\RaggedRight\arraybackslash\columncolor{CornellReferenceCueGray}}p{0.28\textwidth}>{\RaggedRight\arraybackslash}p{0.64\textwidth}@{}}
\rowcolor{CornellReferenceNavy}\color{white}\textbf{Cue / recall prompt} & \color{white}\textbf{Notes}\\\endfirsthead
\rowcolor{CornellReferenceNavy}\color{white}\textbf{Cue / recall prompt} & \color{white}\textbf{Notes (continued)}\\\endhead
\arrayrulecolor{CornellReferenceLineGray}
\textbf{What contract governs General?}\\[-1mm]{\scriptsize\CornellClauseRef{23.4.1}} & Frames the terminology, applicability, and relationships needed to interpret General; detailed rules remain in the following subclauses.\\\hline
\textbf{What contract governs Header <string> synopsis?}\\[-1mm]{\scriptsize\CornellClauseRef{23.4.2}} & Catalogs the declarations made available by Header <string> synopsis. The synopsis is an interface map; the detailed specifications supply constraints and behavior.\\\hline
\textbf{What contract governs Class template basic\_\allowbreak{}string?}\\[-1mm]{\scriptsize\CornellClauseRef{23.4.3}} & Specifies text-sequence behavior for Class template basic\_\allowbreak{}string; establish character type, extent, ownership, null termination, encoding assumptions, and invalidation before use.\\\hline
\textbf{What contract governs Non-member functions?}\\[-1mm]{\scriptsize\CornellClauseRef{23.4.4}} & Specifies the facility family Non-member functions. Read its declarations together with mandates, preconditions, effects, results, exceptions, complexity, and lifetime rules.\\\hline
\textbf{When is Numeric conversions permitted?}\\[-1mm]{\scriptsize\CornellClauseRef{23.4.5}} & Specifies source and destination conditions for Numeric conversions, the resulting type or value category, and any information-loss, pointer, qualification, or lifetime consequence.\\\hline
\textbf{What contract governs Hash support?}\\[-1mm]{\scriptsize\CornellClauseRef{23.4.6}} & Specifies the facility family Hash support. Read its declarations together with mandates, preconditions, effects, results, exceptions, complexity, and lifetime rules.\\\hline
\textbf{What contract governs Suffix for basic\_\allowbreak{}string literals?}\\[-1mm]{\scriptsize\CornellClauseRef{23.4.7}} & Specifies text-sequence behavior for Suffix for basic\_\allowbreak{}string literals; establish character type, extent, ownership, null termination, encoding assumptions, and invalidation before use.\\\hline
\end{longtable}
\begin{CornellReferenceSummaryBox}String classes supplies the library semantics portion of this clause. The key review move is to connect its local wording to the clause-wide model, then classify each condition by normative force before relying on it.\end{CornellReferenceSummaryBox}
\Needspace{5\baselineskip}\paragraph{Detailed coverage ledger.}
\begin{description}[style=nextline,leftmargin=2.8cm,font=\normalfont\bfseries\color{CornellReferenceTeal}]
\item[23.4.3.1] \textbf{General.} Frames the terminology, applicability, and relationships needed to interpret General; detailed rules remain in the following subclauses.
\item[23.4.3.2] \textbf{General requirements.} States the admissibility and semantic obligations for General requirements. Separate compile-time participation rules from caller preconditions and runtime effects.
\item[23.4.3.3] \textbf{Constructors and assignment operators.} Specifies how Constructors and assignment operators establishes object state, including participation constraints, initialization effects, and exceptional exits.
\item[23.4.3.4] \textbf{Iterator support.} Explains the traversal or view contract for Iterator support; prove domain validity, lifetime, sentinel compatibility, and any borrowing or invalidation property.
\item[23.4.3.5] \textbf{Capacity.} Specifies the facility family Capacity. Read its declarations together with mandates, preconditions, effects, results, exceptions, complexity, and lifetime rules.
\item[23.4.3.6] \textbf{Element access.} Specifies the facility family Element access. Read its declarations together with mandates, preconditions, effects, results, exceptions, complexity, and lifetime rules.
\item[23.4.3.7] \textbf{Modifiers.} Specifies the facility family Modifiers. Read its declarations together with mandates, preconditions, effects, results, exceptions, complexity, and lifetime rules.
\item[23.4.3.7.1] \textbf{basic\_\allowbreak{}string::operator+=.} Specifies text-sequence behavior for basic\_\allowbreak{}string::operator+=; establish character type, extent, ownership, null termination, encoding assumptions, and invalidation before use.
\item[23.4.3.7.2] \textbf{basic\_\allowbreak{}string::append.} Specifies text-sequence behavior for basic\_\allowbreak{}string::append; establish character type, extent, ownership, null termination, encoding assumptions, and invalidation before use.
\item[23.4.3.7.3] \textbf{basic\_\allowbreak{}string::assign.} Specifies text-sequence behavior for basic\_\allowbreak{}string::assign; establish character type, extent, ownership, null termination, encoding assumptions, and invalidation before use.
\item[23.4.3.7.4] \textbf{basic\_\allowbreak{}string::insert.} Specifies text-sequence behavior for basic\_\allowbreak{}string::insert; establish character type, extent, ownership, null termination, encoding assumptions, and invalidation before use.
\item[23.4.3.7.5] \textbf{basic\_\allowbreak{}string::erase.} Specifies text-sequence behavior for basic\_\allowbreak{}string::erase; establish character type, extent, ownership, null termination, encoding assumptions, and invalidation before use.
\item[23.4.3.7.6] \textbf{basic\_\allowbreak{}string::replace.} Specifies text-sequence behavior for basic\_\allowbreak{}string::replace; establish character type, extent, ownership, null termination, encoding assumptions, and invalidation before use.
\item[23.4.3.7.7] \textbf{basic\_\allowbreak{}string::copy.} Specifies text-sequence behavior for basic\_\allowbreak{}string::copy; establish character type, extent, ownership, null termination, encoding assumptions, and invalidation before use.
\item[23.4.3.7.8] \textbf{basic\_\allowbreak{}string::swap.} Governs state transfer or exchange for basic\_\allowbreak{}string::swap; check self-operation, validity after movement, exception behavior, and invalidation.
\item[23.4.3.8] \textbf{String operations.} Defines the observable result of String operations together with applicable iterator-domain, ordering, complexity, invalidation, and exception conditions.
\item[23.4.3.8.1] \textbf{Accessors.} Specifies the facility family Accessors. Read its declarations together with mandates, preconditions, effects, results, exceptions, complexity, and lifetime rules.
\item[23.4.3.8.2] \textbf{Searching.} Specifies the facility family Searching. Read its declarations together with mandates, preconditions, effects, results, exceptions, complexity, and lifetime rules.
\item[23.4.3.8.3] \textbf{basic\_\allowbreak{}string::substr.} Specifies text-sequence behavior for basic\_\allowbreak{}string::substr; establish character type, extent, ownership, null termination, encoding assumptions, and invalidation before use.
\item[23.4.3.8.4] \textbf{basic\_\allowbreak{}string::compare.} Specifies text-sequence behavior for basic\_\allowbreak{}string::compare; establish character type, extent, ownership, null termination, encoding assumptions, and invalidation before use.
\item[23.4.3.8.5] \textbf{basic\_\allowbreak{}string::starts\_\allowbreak{}with.} Specifies text-sequence behavior for basic\_\allowbreak{}string::starts\_\allowbreak{}with; establish character type, extent, ownership, null termination, encoding assumptions, and invalidation before use.
\item[23.4.3.8.6] \textbf{basic\_\allowbreak{}string::ends\_\allowbreak{}with.} Specifies text-sequence behavior for basic\_\allowbreak{}string::ends\_\allowbreak{}with; establish character type, extent, ownership, null termination, encoding assumptions, and invalidation before use.
\item[23.4.3.8.7] \textbf{basic\_\allowbreak{}string::contains.} Specifies text-sequence behavior for basic\_\allowbreak{}string::contains; establish character type, extent, ownership, null termination, encoding assumptions, and invalidation before use.
\item[23.4.4.1] \textbf{operator+.} Specifies the facility family operator+. Read its declarations together with mandates, preconditions, effects, results, exceptions, complexity, and lifetime rules.
\item[23.4.4.2] \textbf{Non-member comparison operator functions.} Specifies the facility family Non-member comparison operator functions. Read its declarations together with mandates, preconditions, effects, results, exceptions, complexity, and lifetime rules.
\item[23.4.4.3] \textbf{swap.} Governs state transfer or exchange for swap; check self-operation, validity after movement, exception behavior, and invalidation.
\item[23.4.4.4] \textbf{Inserters and extractors.} Specifies the facility family Inserters and extractors. Read its declarations together with mandates, preconditions, effects, results, exceptions, complexity, and lifetime rules.
\item[23.4.4.5] \textbf{Erasure.} Specifies the facility family Erasure. Read its declarations together with mandates, preconditions, effects, results, exceptions, complexity, and lifetime rules.
\end{description}
\Needspace{8\baselineskip}
\subsection*{23.5 Null-terminated sequence utilities}
\begin{longtable}{@{}>{\RaggedRight\arraybackslash\columncolor{CornellReferenceCueGray}}p{0.28\textwidth}>{\RaggedRight\arraybackslash}p{0.64\textwidth}@{}}
\rowcolor{CornellReferenceNavy}\color{white}\textbf{Cue / recall prompt} & \color{white}\textbf{Notes}\\\endfirsthead
\rowcolor{CornellReferenceNavy}\color{white}\textbf{Cue / recall prompt} & \color{white}\textbf{Notes (continued)}\\\endhead
\arrayrulecolor{CornellReferenceLineGray}
\textbf{What contract governs Header <cctype> synopsis?}\\[-1mm]{\scriptsize\CornellClauseRef{23.5.1}} & Catalogs the declarations made available by Header <cctype> synopsis. The synopsis is an interface map; the detailed specifications supply constraints and behavior.\\\hline
\textbf{What contract governs Header <cwctype> synopsis?}\\[-1mm]{\scriptsize\CornellClauseRef{23.5.2}} & Catalogs the declarations made available by Header <cwctype> synopsis. The synopsis is an interface map; the detailed specifications supply constraints and behavior.\\\hline
\textbf{What contract governs Header <cstring> synopsis?}\\[-1mm]{\scriptsize\CornellClauseRef{23.5.3}} & Catalogs the declarations made available by Header <cstring> synopsis. The synopsis is an interface map; the detailed specifications supply constraints and behavior.\\\hline
\textbf{What contract governs Header <cwchar> synopsis?}\\[-1mm]{\scriptsize\CornellClauseRef{23.5.4}} & Catalogs the declarations made available by Header <cwchar> synopsis. The synopsis is an interface map; the detailed specifications supply constraints and behavior.\\\hline
\textbf{What contract governs Header <cuchar> synopsis?}\\[-1mm]{\scriptsize\CornellClauseRef{23.5.5}} & Catalogs the declarations made available by Header <cuchar> synopsis. The synopsis is an interface map; the detailed specifications supply constraints and behavior.\\\hline
\textbf{When is Multibyte / wide string and character conversion functions permitted?}\\[-1mm]{\scriptsize\CornellClauseRef{23.5.6}} & Specifies source and destination conditions for Multibyte / wide string and character conversion functions, the resulting type or value category, and any information-loss, pointer, qualification, or lifetime consequence.\\\hline
\end{longtable}
\begin{CornellReferenceSummaryBox}Null-terminated sequence utilities supplies the library semantics portion of this clause. The key review move is to connect its local wording to the clause-wide model, then classify each condition by normative force before relying on it.\end{CornellReferenceSummaryBox}
\section{Language-Lawyer and Library-Usage Pitfalls}
\begin{CornellReferencePitfallBox}\begin{itemize}[label=\textcolor{CornellReferenceAmber}{$\blacktriangleright$}]
\item Reading a synopsis as the complete behavioral contract.
\item Ignoring mandates, preconditions, exception behavior, or complexity requirements.
\item Using an exposition-only name as though it were public API.
\item Assuming invalidation, ownership, or thread-safety properties that are not stated.
\item Treating successful compilation as proof that semantic requirements and preconditions are met.
\end{itemize}\end{CornellReferencePitfallBox}
\section{Programmer and Implementation Checklist}
\begin{CornellReferenceChecklistBox}\begin{itemize}[label=$\square$]
\item Identify the declaring header or module exposure for each facility.
\item Check template arguments and mandates before evaluating an operation.
\item Establish every precondition at the call site.
\item Record effects, returned value, postconditions, and exceptions separately.
\item Audit iterator, reference, pointer, and object lifetime after mutation.
\item Verify stated complexity and synchronization assumptions.
\item For \CornellClauseRef{23.1}, verify general against the precise rule category and all cited dependencies.
\item For \CornellClauseRef{23.2}, verify character traits against the precise rule category and all cited dependencies.
\item For \CornellClauseRef{23.3}, verify string view classes against the precise rule category and all cited dependencies.
\item For \CornellClauseRef{23.4}, verify string classes against the precise rule category and all cited dependencies.
\end{itemize}\end{CornellReferenceChecklistBox}
\section{Clause Synthesis}
\begin{CornellReferenceOrientationBox}Defines character traits, owning strings, string views, null-terminated sequence utilities, numeric conversions, and related hashing and I/O. It distinguishes ownership, bounds, encoding units, and invalidation behavior. The primary character-sequence abstraction layer. A sound review distinguishes syntax and participation from runtime preconditions, keeps notes and examples non-mandatory, and follows cross-clause dependencies before making a portability claim.\end{CornellReferenceOrientationBox}
\section{Self-Test Questions}
\begin{enumerate}
\item What role does General play in the clause's overall model?
\item Which normative distinctions must be kept separate when applying Character traits?
\item How would you audit a program or implementation that depends on String view classes?
\item Compare String classes with Null-terminated sequence utilities; what different question does each answer?
\item What role does Null-terminated sequence utilities play in the clause's overall model?
\item Which normative distinctions must be kept separate when applying General?
\item How would you audit a program or implementation that depends on Character traits requirements?
\item Compare Traits typedefs with char\_\allowbreak{}traits specializations; what different question does each answer?
\item What role does char\_\allowbreak{}traits specializations play in the clause's overall model?
\item Which normative distinctions must be kept separate when applying General?
\item How would you audit a program or implementation that depends on Header <string\_\allowbreak{}view> synopsis?
\item Compare Class template basic\_\allowbreak{}string\_\allowbreak{}view with Non-member comparison functions; what different question does each answer?
\item What role does Non-member comparison functions play in the clause's overall model?
\item Which normative distinctions must be kept separate when applying Inserters and extractors?
\item Scenario: a program relies on an unstated implementation behavior while using Hash support. What must be checked before calling the program portable?
\item Scenario: a program relies on an unstated implementation behavior while using Suffix for basic\_\allowbreak{}string\_\allowbreak{}view literals. What must be checked before calling the program portable?
\end{enumerate}
\clearpage\section{Answer Key}
\begin{enumerate}
\item General is one part of the clause's library semantics model. Its role is understood by connecting the local rule in 23.4.3.1 to the clause orientation and its dependent concepts.
\item Keep formation and well-formedness separate from runtime preconditions and effects. Also distinguish mandatory wording from notes, implementation choice from undefined behavior, and interface declarations from detailed contracts; see 23.2.
\item Audit String view classes by locating the controlling wording in 23.3, listing inputs and type constraints, proving any preconditions, recording effects and failure behavior, and checking lifetime, invalidation, complexity, and synchronization where applicable.
\item String classes and Null-terminated sequence utilities occupy different steps or facility groups. Analyze each under its own subclause, then follow the explicit dependency between them rather than transferring one set of guarantees to the other; begin at 23.4.
\item Null-terminated sequence utilities is one part of the clause's library semantics model. Its role is understood by connecting the local rule in 23.5 to the clause orientation and its dependent concepts.
\item Keep formation and well-formedness separate from runtime preconditions and effects. Also distinguish mandatory wording from notes, implementation choice from undefined behavior, and interface declarations from detailed contracts; see 23.4.3.1.
\item Audit Character traits requirements by locating the controlling wording in 23.2.2, listing inputs and type constraints, proving any preconditions, recording effects and failure behavior, and checking lifetime, invalidation, complexity, and synchronization where applicable.
\item Traits typedefs and char\_\allowbreak{}traits specializations occupy different steps or facility groups. Analyze each under its own subclause, then follow the explicit dependency between them rather than transferring one set of guarantees to the other; begin at 23.2.3.
\item char\_\allowbreak{}traits specializations is one part of the clause's library semantics model. Its role is understood by connecting the local rule in 23.2.4 to the clause orientation and its dependent concepts.
\item Keep formation and well-formedness separate from runtime preconditions and effects. Also distinguish mandatory wording from notes, implementation choice from undefined behavior, and interface declarations from detailed contracts; see 23.4.3.1.
\item Audit Header <string\_\allowbreak{}view> synopsis by locating the controlling wording in 23.3.2, listing inputs and type constraints, proving any preconditions, recording effects and failure behavior, and checking lifetime, invalidation, complexity, and synchronization where applicable.
\item Class template basic\_\allowbreak{}string\_\allowbreak{}view and Non-member comparison functions occupy different steps or facility groups. Analyze each under its own subclause, then follow the explicit dependency between them rather than transferring one set of guarantees to the other; begin at 23.3.3.
\item Non-member comparison functions is one part of the clause's object contract model. Its role is understood by connecting the local rule in 23.3.4 to the clause orientation and its dependent concepts.
\item Keep formation and well-formedness separate from runtime preconditions and effects. Also distinguish mandatory wording from notes, implementation choice from undefined behavior, and interface declarations from detailed contracts; see 23.4.4.4.
\item First identify the exact requirement in 23.4.6, then classify the relied-on behavior as required, unspecified, implementation-defined, conditionally supported, or undefined. Portability requires satisfying the program obligations and avoiding dependence on a choice not guaranteed in the target environments.
\item First identify the exact requirement in 23.3.7, then classify the relied-on behavior as required, unspecified, implementation-defined, conditionally supported, or undefined. Portability requires satisfying the program obligations and avoiding dependence on a choice not guaranteed in the target environments.
\end{enumerate}
\section{Key-Term Recap}
\begin{longtable}{@{}>{\RaggedRight\arraybackslash}p{0.28\textwidth}>{\RaggedRight\arraybackslash}p{0.64\textwidth}@{}}\toprule\textbf{Term} & \textbf{Working meaning}\\\midrule\endfirsthead\toprule\textbf{Term} & \textbf{Working meaning (continued)}\\\midrule\endhead\bottomrule\endfoot
char\_\allowbreak{}traits & An operations policy for a character type.\\
basic\_\allowbreak{}string & An owning contiguous character container.\\
basic\_\allowbreak{}string\_\allowbreak{}view & A non-owning view over a character sequence.\\
NTCTS & A sequence ending at a null character value.\\
npos & A sentinel position value used by string and view operations.\\
to\_\allowbreak{}chars & A locale-independent value-to-character conversion family.\\
\end{longtable}
\CornellTakeaway{String correctness starts by separating ownership, extent, character units, encoding assumptions, and null termination.}
\end{document}
